\subsection{The $\{\minFac, \operatorname{secondMinFac}\}$ Variant}
\label{sec:minfac-secondminfac}

The two-point variant interpolates between the random and deterministic
extremes.  At each step, instead of a single deterministic choice or a
fully random one, two specific factors are available:
$\minFac(\Prod{2}(n)+1)$ and $\operatorname{secondMinFac}(\Prod{2}(n)+1)$.
A decision sequence $\sigma : \NN \to \mathrm{Bool}$ selects between
them; the all-true sequence recovers the standard EM walk.

\paragraph{The $\varepsilon$-walk.}
The formalization defines a
\textbf{self-consistent $\varepsilon$-walk}
(\lean{EM/Stochastic/VanishingNoiseVariant.lean}{epsWalkProd})
in which the chosen factor---either $\minFac(P(n){+}1)$ or
$\operatorname{secondMinFac}(P(n){+}1)$---becomes the actual multiplier,
updating the accumulator and producing a new Euclid number at the next
step.
The function $\operatorname{secondMinFac}(n)$ (second-smallest prime
factor) is defined and proved prime, divisor of~$n$, and
$\geq \minFac(n)$
(\lean{EM/Stochastic/VanishingNoiseVariant.lean}{secondMinFac\_prime},
\lean{EM/Stochastic/VanishingNoiseVariant.lean}{secondMinFac\_dvd},
\lean{EM/Stochastic/VanishingNoiseVariant.lean}{minFac\_le\_secondMinFac}).

This self-consistent walk generates a \emph{binary branching tree}
of depth~$N$ with $2^N$ leaves, where each internal node has its own
path-dependent factor set.  The \emph{tree character sum}
(\lean{EM/Stochastic/VanishingNoiseVariant.lean}{treeCharSum}),
a weighted sum of $\chi(\text{leaf position})$ over all leaves, is proved
$\leq 1$ in norm.  Unlike the product multiset of the stochastic MC, this sum
does \textbf{not} factorize as $\prod_k \bar\chi(S_k)$: the left and
right subtrees start from different accumulators, so their character sums
differ.  The open hypothesis
\defname{TreeContractionHypothesis}
asserts that the tree character sum tends to zero for all nontrivial~$\chi$;
it is implied by CME (strictness unknown) and strictly stronger than the product approximation.

\paragraph{Non-self-consistent variant MC.}
An independent route bypasses the tree contraction problem entirely.
At each EM step~$n$, form the two-element unit set
$\{\minFac(P(n){+}1),\, \operatorname{secondMinFac}(P(n){+}1)\} \bmod q$
in $(\ZZ/q\ZZ)^\times$.  When one factor equals~$q$
(a non-unit), fall back to $\operatorname{Finset.univ}$, whose
mean character value is~$0$ by orthogonality---giving perfect contraction
(\lean{EM/Stochastic/VanishingNoiseVariantB.lean}{meanCharValue\_univ\_eq\_zero}).
Under \defname{UFDStrong}---spectral gaps of the padded unit sets
are non-summable for every nontrivial~$\chi$---the
sparse product contraction machinery gives vanishing averaged
products, and the Fourier counting argument yields
path existence.  Concretely:

\begin{theorem}[\lean{EM/Stochastic/VanishingNoiseVariantB.lean}{variant\_mc\_from\_ufd\_strong\_proved}]
\label{thm:variant-mc}
For every prime $q \geq 3$, if $\mathrm{UFDStrong}(q)$ holds, then
every element of $(\ZZ/q\ZZ)^\times$ is representable as a product
of selections from $\{\minFac,\operatorname{secondMinFac}\}$ of
consecutive Euclid numbers.  In particular, some selection path
hits~$-1$.
\end{theorem}

The verified chain is:
\[
  \mathrm{UFDStrong}(q)
  \;\Longrightarrow\;
  \text{vanishing products}
  \;\Longrightarrow\;
  \text{path existence}
  \;\Longrightarrow\;
  \text{VariantMC}.
\]
Every arrow is machine-checked with zero \code{sorry}.

\paragraph{The gap: UFDStrong.}
\defname{UFDStrong}$(q)$ asserts: for each nontrivial character
$\chi$ on $(\ZZ/q\ZZ)^\times$, the spectral gaps
$1 - \|\bar\chi(\text{paddedUnitSet}_n)\|$ are non-summable.
The hypothesis splits into two cases:
\begin{enumerate}[nosep]
\item \emph{Fallback case}: if infinitely many steps have
  $\minFac \equiv \operatorname{secondMinFac} \pmod{q}$ (or one
  equals~$q$), the padded set falls back to the full group, giving
  perfect contraction (gap~$= 1$) at each such step.
  Non-summability is immediate.
  \textbf{Proved}:
  \lean{EM/Stochastic/VanishingNoiseVariantB.lean}{ufd\_fallback\_not\_summable}.
\item \emph{Diversity case}: if cofinitely many steps have two
  distinct units mod~$q$, one needs the character values at those
  units to differ infinitely often---i.e., the ratio
  $\minFac/\operatorname{secondMinFac}$ does not persistently lie in
  $\ker(\chi)$.  This is a genuine number-theoretic question about
  Euclid number factorizations.  \textbf{Open}:
  \code{UFDImpliesUFDStrong}.
\end{enumerate}

\paragraph{Variant MC does not imply standard MC.}
The selection~$\sigma$ that achieves hitting may differ per
prime~$q$, and the virtual accumulator
$\prod_k \text{factor}_{\sigma(k)}$ is not the EM accumulator
$P(n)$.  The gap is the same orbit-specificity barrier
(Dead End~\#130,
\lean{EM/Group/CyclicWalkCoverage.lean}{alternating\_walk\_misses\_two}):
generation does not imply coverage for deterministic walks.

\paragraph{Three routes to UFDStrong.}
The formalization identifies a hierarchy of progressively weaker
sufficient conditions for UFDStrong$(q)$:

\begin{enumerate}[nosep]
\item \textbf{Quantitative escape} (\defname{MinFacRatioEscape}):
  for each nontrivial~$\chi$, there exists $\delta > 0$ such that
  infinitely many steps have spectral gap $\geq \delta$.
  Since a non-summable sequence with a positive lower bound infinitely
  often is non-summable, this implies UFDStrong directly
  (\lean{EM/Stochastic/VanishingNoiseVariantB.lean}{ratio\_escape\_implies\_ufdStrong}).

\item \textbf{Qualitative escape} (\defname{MinFacRatioEscapeQual}):
  for each nontrivial~$\chi$, infinitely many steps have
  $\operatorname{card}(\text{rawTwoPointUnitSet}) \geq 2$ \emph{and}
  $\chi(\minFac) \neq \chi(\operatorname{secondMinFac})$.
  This implies quantitative escape via a finite-group pigeonhole
  argument: on a finite group, the gap function takes finitely many
  values, so being positive infinitely often gives a uniform positive
  lower bound
  (\lean{EM/Stochastic/VanishingNoiseVariantB.lean}{qual\_implies\_quant}).

\item \textbf{Orbit-level minFac equidistribution} (\defname{OrbitMFRE}):
  the residues $\minFac(\Prod{2}(n){+}1) \bmod q$ are equidistributed
  among the units of $(\ZZ/q\ZZ)^\times$ in Cesàro density.
  This is the weakest condition in the hierarchy: equidistributed
  residues prevent persistent chi-degeneracy.  The bridge
  OrbitMFRE $\Rightarrow$ MinFacRatioEscapeQual is \textbf{open}
  (\code{OrbitMFREImpliesEscapeQual}).
\end{enumerate}

The assembled chain is:
\[
  \underbrace{\mathrm{OrbitMFRE}}_{\text{open}}
  \;\xRightarrow{\text{open bridge}}\;
  \mathrm{EscapeQual}
  \;\xRightarrow{\text{proved}}\;
  \mathrm{EscapeQuant}
  \;\xRightarrow{\text{proved}}\;
  \mathrm{UFDStrong}
  \;\xRightarrow{\text{proved}}\;
  \mathrm{VariantMC}.
\]
Every arrow except the first is machine-checked with zero \code{sorry}.
The most concrete attack target is MinFacRatioEscapeQual: one needs to
show that $\minFac$ and $\operatorname{secondMinFac}$ do not
persistently agree modulo every character.  OrbitMFRE faces the same
orbit-specificity barrier as CME for standard MC.

\paragraph{Stochastic Two-Point MC.}
The non-self-consistent machinery yields a clean quantitative statement.
Define the \defname{StochasticTwoPointMC} problem: under UFDStrong$(q)$
with $q \geq 3$, does every element of $(\ZZ/q\ZZ)^\times$ appear with
positive count in the product multiset of padded unit sets?

\begin{theorem}[\lean{EM/Stochastic/VanishingNoiseVariantC.lean}{stochastic\_two\_point\_mc\_proved}]
\label{thm:stochastic-two-point}
$\mathrm{UFDStrong}(q) \;\Longrightarrow\; \mathrm{StochasticTwoPointMC}(q)$.
\end{theorem}

\noindent The product multiset grows exponentially
($\geq 2^N$ paths for $q \geq 3$,
\lean{EM/Stochastic/VanishingNoiseVariantC.lean}{productMultiset\_card\_ge\_two\_pow}).
A \defname{fairCoin} schedule ($\varepsilon = 1/2$ at every step) gives
the \defname{fairTreeCharSum}---the tree character sum at constant
$1/2$---which is bounded by~$1$ in norm
(\lean{EM/Stochastic/VanishingNoiseVariantC.lean}{fairTreeCharSum\_norm\_le\_one}).
TreeContractionHypothesis implies the weaker
\defname{TreeContractionAtHalf}
(\lean{EM/Stochastic/VanishingNoiseVariantC.lean}{treeContraction\_of\_hypothesis}),
and at $\varepsilon = 0$ the tree degenerates to the deterministic
$\minFac$ path
(\lean{EM/Stochastic/VanishingNoiseVariantC.lean}{treeCharSum\_at\_zero\_step}).

The gap between the non-self-consistent result (proved) and a
\emph{self-consistent} stochastic MC is the \textbf{selection coherence}
problem: the non-self-consistent product multiset uses factor sets from
the actual EM walk (which always picks $\minFac$), so the factor sets
are fixed and the product factorizes; in the self-consistent tree,
picking $\operatorname{secondMinFac}$ at step~$n$ changes the
accumulator, hence the factor set at step~$n{+}1$, preventing
factorization.

\paragraph{Faithful character escape.}
A character $\chi$ on $(\ZZ/q\ZZ)^\times$ is \defname{faithful} if it
is injective (equivalently, $\ker(\chi) = \{1\}$).  For faithful
characters, whenever the raw two-point unit set has two distinct
elements, their $\chi$-values are automatically distinct---giving a
positive spectral gap at every non-fallback step, with no
number-theoretic hypothesis.  Combined with the fallback case (gap $= 1$
at fallback steps, already proved), this yields unconditional
non-summability:

\begin{theorem}[\lean{EM/Stochastic/VanishingNoiseVariantC.lean}{faithful\_character\_escape}]
\label{thm:faithful-escape}
For every prime $q \geq 3$ and every faithful nontrivial character
$\chi$ on $(\ZZ/q\ZZ)^\times$:
\[
  \sum_n \bigl(1 - \|\bar\chi(\text{paddedUnitSet}_n)\|\bigr)
  = +\infty.
\]
Unconditional---zero open hypotheses.
\end{theorem}

\noindent For groups of prime order, every nontrivial character is
faithful (by Lagrange: the kernel is a subgroup of a prime-order group,
hence trivial or the whole group;
\lean{EM/Stochastic/VanishingNoiseVariantC.lean}{prime\_order\_nontrivial\_is\_faithful}).
Since $|(\ZZ/3\ZZ)^\times| = 2$ is prime, UFDStrong$(3)$ holds
\emph{unconditionally}
(\lean{EM/Stochastic/VanishingNoiseVariantC.lean}{ufdStrong\_three}), and
the full chain closes:

\begin{theorem}[\lean{EM/Stochastic/VanishingNoiseVariantC.lean}{variant\_mc\_three\_unconditional}]
\label{thm:variant-mc-3}
$\mathrm{StochasticTwoPointMC}(3)$ holds with zero open hypotheses:
every element of $(\ZZ/3\ZZ)^\times$ appears with positive count
in the product multiset.
\end{theorem}

\noindent For primes $q \geq 5$ where $|(\ZZ/q\ZZ)^\times| = q{-}1$ is
composite, non-faithful nontrivial characters exist, and the sole
remaining open hypothesis is
\defname{NonFaithfulCharacterEscape}---spectral gap non-summability
for characters with nontrivial kernel.  The faithful/non-faithful
decomposition is clean:

\begin{theorem}[\lean{EM/Stochastic/VanishingNoiseVariantC.lean}{nfce\_implies\_ufdStrong}]
\label{thm:nfce-ufdstrong}
$\mathrm{NonFaithfulCharacterEscape}(q)
\;\Longrightarrow\; \mathrm{UFDStrong}(q)$
for all primes $q \geq 3$.
\end{theorem}

\noindent The faithful half is free; only the non-faithful half
remains open.

\paragraph{NFCE infrastructure.}
The formalization develops a precise characterization of what NFCE
requires.  Define the \defname{factorRatio} at step~$n$ as
$\mathrm{liftToUnit}(\minFac) \cdot
\mathrm{liftToUnit}(\operatorname{secondMinFac})^{-1}$
in~$(\ZZ/q\ZZ)^\times$.
At non-fallback steps (raw card $\geq 2$), the spectral gap is zero
if and only if this ratio lies in~$\ker(\chi)$
(\lean{EM/Stochastic/NonFaithfulCharacterEscape.lean}{gap\_zero\_iff\_ratio\_in\_ker}).
NFCE thus reduces to: for each non-faithful nontrivial~$\chi$, the
factor ratio does not eventually stay in~$\ker(\chi)$.

If NFCE fails for a specific~$\chi$, the gaps are summable,
which forces \defname{KernelConfinement}: the ratio is eventually
confined to~$\ker(\chi)$
(\lean{EM/Stochastic/NonFaithfulCharacterEscape.lean}{summable\_implies\_ratio\_confined}).
Since $\ker(\chi)$ has index $\geq 2$ for nontrivial~$\chi$
(\lean{EM/Stochastic/NonFaithfulCharacterEscape.lean}{ker\_index\_ge\_two}),
kernel confinement is a strong arithmetic constraint:
$\minFac / \operatorname{secondMinFac} \bmod q$ eventually lies in a
proper subgroup of the unit group.

\paragraph{NFCE(5): reduction to a single quadratic escape.}
For $q = 5$, the unit group $(\ZZ/5\ZZ)^\times$ has order~$4$
(cyclic, $\cong \ZZ/4\ZZ$).  Its characters partition as: the trivial
character, two faithful characters (order~$4$, handled unconditionally
by FaithfulCharacterEscape), and one non-faithful nontrivial character
(order~$2$, the unique quadratic character).  Thus NFCE(5) reduces to a
single escape condition: the factor ratio
$\minFac / \operatorname{secondMinFac} \bmod 5$ must leave the unique
index-2 subgroup of $(\ZZ/5\ZZ)^\times$ infinitely often.
The formalization proves that this
\defname{RatioKernelEscape} condition suffices for the full chain
\lean{EM/Stochastic/NonFaithfulCharacterEscape.lean}{variant\_mc\_five\_of\_ratio\_escape}
{\color{leanink}$\checkmark$}:
RatioKernelEscape $\Rightarrow$ NFCE(5) $\Rightarrow$ UFDStrong(5)
$\Rightarrow$ StochasticTwoPointMC(5).

\paragraph{Intersection kernel dichotomy.}
When NFCE fails for \emph{all} non-faithful nontrivial characters
simultaneously (\defname{TotalNFCEFailure}), the factor ratio is
eventually confined to every non-faithful kernel (from the
summability-to-confinement lemma of Part~25).
The intersection of these kernels is constrained by a group-theoretic
fact: in a finite group whose order has $\geq 2$ distinct prime
factors, there exist subgroups of coprime orders whose intersection is
trivial
(\lean{EM/Stochastic/NonFaithfulCharacterEscape.lean}{coprime\_order\_zpowers\_inf\_bot}
{\color{leanink}$\checkmark$}).  Combined with Cauchy's theorem, this
yields \defname{NonFaithfulCharSeparation}~(NFCS): for every
$g \neq 1$, some non-faithful nontrivial character has $\chi(g) \neq 1$.
NFCS is
\textbf{proved} for all groups whose order has $\geq 2$ distinct prime
factors
(\lean{EM/Stochastic/NonFaithfulCharacterEscape.lean}{nonFaithfulCharSeparation\_of\_two\_prime\_factors}
{\color{leanink}$\checkmark$}), which for $(\ZZ/q\ZZ)^\times$ covers
all primes except $q = 2, 3, 5$ and Fermat primes (where $q - 1$ is a
power of~$2$).
Under NFCS, TotalNFCEFailure forces
$\minFac / \operatorname{secondMinFac} \equiv 1 \pmod{q}$ eventually
(\lean{EM/Stochastic/NonFaithfulCharacterEscape.lean}{total\_confinement\_implies\_ratio\_one}
{\color{leanink}$\checkmark$}).
This yields:

\begin{theorem}[Intersection kernel dichotomy]
\label{thm:kernel-dichotomy}
Let $q$ be prime with $q - 1$ having $\geq 2$ distinct prime factors.
Then exactly one of the following holds:
\begin{enumerate}[nosep]
\item $\mathrm{NFCE}(q)$ holds, giving
  $\mathrm{UFDStrong}(q)$ and $\mathrm{StochasticTwoPointMC}(q)$; or
\item $\minFac(\Prod{2}(n)+1) \equiv \operatorname{secondMinFac}(\Prod{2}(n)+1) \pmod{q}$
  for all sufficiently large~$n$.
\end{enumerate}
The second alternative is a strong arithmetic rigidity: the two
smallest prime factors of each Euclid number are eventually
indistinguishable mod~$q$.
\end{theorem}

\noindent
The proof combines NonFaithfulCharSeparation
(\lean{EM/Stochastic/NonFaithfulCharacterEscape.lean}{nonFaithfulCharSeparation\_of\_two\_prime\_factors}
{\color{leanink}$\checkmark$})
with the summability-to-confinement lemma
(\lean{EM/Stochastic/NonFaithfulCharacterEscape.lean}{total\_confinement\_implies\_ratio\_one}
{\color{leanink}$\checkmark$}).
The condition ``$q{-}1$ has $\geq 2$ distinct prime factors'' covers
all primes except $q = 2, 3, 5$ and Fermat primes (where $q{-}1$ is a
power of~$2$).

\paragraph{Quotient character lift.}
The quotient character infrastructure, formalized in Part~28 of
\code{EM/Stochastic/NonFaithfulCharacterEscape.lean} (173~lines),
provides a structural reduction of NFCE\@.
For a non-faithful nontrivial character~$\chi$, the kernel
$\ker(\chi)$ is a proper nontrivial subgroup; by the universal
property, $\chi$ factors as
$\chi = \bar\chi \circ \pi$
where $\pi \colon G \to G/\ker(\chi)$ is the quotient map.
The induced character~$\bar\chi$ on the quotient group is
\textbf{faithful} (injective)
(\lean{EM/Stochastic/NonFaithfulCharacterEscape.lean}{quotientChar\_faithful}
{\color{leanink}$\checkmark$}),
the quotient has order $\geq 2$
(\lean{EM/Stochastic/NonFaithfulCharacterEscape.lean}{quotient\_card\_ge\_two}
{\color{leanink}$\checkmark$}),
and KernelConfinement is equivalent to the quotient projection
of the factor ratio being eventually trivial
(\lean{EM/Stochastic/NonFaithfulCharacterEscape.lean}{kernelConfinement\_iff\_quotient\_eventually\_one}
{\color{leanink}$\checkmark$}).
\begin{theorem}[Quotient character lift]
\label{thm:quotient-lift}
For any non-faithful nontrivial character~$\chi$ on~$(\ZZ/q\ZZ)^\times$:
\begin{enumerate}[nosep]
\item $\chi = \bar\chi \circ \pi$ where $\bar\chi$ is faithful
  on $G/\ker(\chi)$ and $|G/\ker(\chi)| \geq 2$.
\item $\mathrm{KernelConfinement}(\chi) \;\Longleftrightarrow\;$
  the quotient projection of the factor ratio is eventually trivial.
\item $\mathrm{RatioKernelEscape}(\chi) \;\Longleftrightarrow\;$
  the quotient projection is nontrivial infinitely often.
\end{enumerate}
\end{theorem}

\noindent This reduces NFCE to: the image of factorRatio in each quotient
$G/\ker(\chi)$ escapes~$\{1\}$ infinitely often.
Since $\bar\chi$ is faithful on the (smaller) quotient group, the
unconditional FaithfulCharacterEscape result
(Theorem~\ref{thm:faithful-escape}) applies to $G/\ker(\chi)$,
connecting the non-faithful case to the faithful framework.

\paragraph{Comparison of open hypotheses.}
Standard MC requires CME (orbit-level conditional equidistribution of
$\minFac$ residues).  The $\{\minFac, \operatorname{secondMinFac}\}$ variant
requires UFDStrong (spectral diversity of two-point factor sets), and
for $q = 3$ this is now \textbf{unconditional}; for $q \geq 5$ the
residual hypothesis is NonFaithfulCharacterEscape.
Neither CME nor UFDStrong implies the other: CME is about the
\emph{selected} factor's equidistribution; UFDStrong is about the
\emph{pair}'s spectral gap.
But both are instances of the same meta-question: does the deterministic
EM construction produce enough diversity in its factorizations to
prevent persistent alignment with any fixed character?

\paragraph{Ensemble two-point: population-level reduction.}
The ensemble framework of Section~\ref{sec:ensemble} can also be
applied to the two-point variant.  For generalized EM starting from
squarefree~$n$, the \defname{genFactorRatio} at step~$k$ is the ratio
$\minFac / \operatorname{secondMinFac} \bmod q$ in
$(\ZZ/q\ZZ)^\times$.  For a nontrivial character~$\chi$, a step
``escapes'' if this ratio lies outside~$\ker(\chi)$---producing a
spectral gap.  The formalization
(\code{EM/Ensemble/TwoPointEnsemble.lean}, 566~lines, zero \code{sorry})
develops the full first-moment chain:
\[
  \text{PRED}
  \;\xrightarrow[\text{Fubini}]{\checkmark}\;
  \text{linear first moment}
  \;\xrightarrow[\text{partition}]{\checkmark}\;
  \text{positive density of high escapers.}
\]
Here \textbf{PopulationRatioEscapeDensity}~(PRED) asserts that for
each nontrivial~$\chi$ and each step $k \geq 1$, a positive
density~$\delta > 0$ of squarefree starting points have factor ratio
outside~$\ker(\chi)$.  Under PRED, the ensemble average of cumulative
escape counts grows linearly
(\lean{EM/Ensemble/TwoPointEnsemble.lean}{escape\_first\_moment\_linear}
{\color{leanink}$\checkmark$}), and the partition argument gives
\lean{EM/Ensemble/TwoPointEnsemble.lean}{positive\_density\_high\_escape}
{\color{leanink}$\checkmark$}: for each threshold~$M$, a positive
density of starting points achieve $\geq M$ cumulative escapes
within~$K$ steps.  PRED itself is open; an earlier draft expected it
from MFRE (MinFacResidueEquidist), which is false (Dead End~\#160), so
the expected route is gone (\code{MFREImpliesPopulationRatioEscape} is
archived).

\paragraph{Self-consistent random walk: the Fourier bridge.}
The stochastic two-point variant above fixes the factor sets
$\{p_1(n), p_2(n)\}$ from the \emph{actual} EM walk and randomizes
the selection.  A stronger variant makes the construction
\emph{self-consistent}: at each step, the accumulator determines the
available factors, and the fair-coin choice feeds back into the next
accumulator.  Formally, define
$\mathrm{epsWalkProdFrom}(\mathrm{acc}, \sigma, 0) = \mathrm{acc}$,
$\mathrm{epsWalkProdFrom}(\mathrm{acc}, \sigma, n{+}1)
= P_n \cdot (\text{if } \sigma(n) \text{ then } \minFac(P_n{+}1)
\text{ else } \operatorname{secondMinFac}(P_n{+}1))$
where $P_n = \mathrm{epsWalkProdFrom}(\mathrm{acc}, \sigma, n)$.
The decision sequence $\sigma : \NN \to \mathrm{Bool}$ replaces the
deterministic $\minFac$ rule with a random binary choice.

The central result is the \textbf{Fourier bridge identity}
(\lean{EM/Stochastic/RandomTwoPointMC.lean}{fourier\_bridge\_identity\_proved}
{\color{leanink}$\checkmark$}): for any accumulator $\mathrm{acc} \geq 2$,
\[
  \mathrm{treeCharSum}(q, \chi, N, \mathrm{acc}, \tfrac{1}{2})
  \;=\;
  \frac{1}{2^N}
  \sum_{\sigma \in \{0,1\}^N}
  \prod_{k < N} \mathrm{chiAt}(q, \chi,
    \mathrm{epsWalkFactorFrom}(\mathrm{acc}, \sigma, k)).
\]
The proof is by induction on~$N$: at each step, the $2^N$ paths split
by first bit into two halves corresponding to the left and right
subtrees of the tree recursion; shift lemmas
(\code{epsWalkProdFrom\_shift}, \code{epsWalkFactorFrom\_shift})
connect the concatenated walk to the shifted walk at the new
accumulator; the $\tfrac{1}{2}$ weights combine to match the tree
recurrence.  Specializing to $\mathrm{acc} = 2$ gives the
unconditional identity
$\mathrm{fairTreeCharSum} = \mathrm{pathCharSum}$
(\lean{EM/Stochastic/RandomTwoPointMC.lean}{fairTreeCharSum\_eq\_pathCharSum}
{\color{leanink}$\checkmark$}).

\defname{RandomTwoPointMC}
(\lean{EM/Stochastic/RandomTwoPointMC.lean}{RandomTwoPointMC})
asks: for $q \geq 3$ and every unit $a \in (\ZZ/q\ZZ)^\times$, does
some binary decision sequence of finite length produce a walk from
$\mathrm{acc} = 2$ landing in residue class~$a$?  This is the
self-consistent analog of StochasticTwoPointMC.
The Fourier bridge connects the tree to the path character sum.
To connect path character sums to reached residue classes, one needs
\defname{ChiAtMultiplicativity}: that $\mathrm{chiAt}(\mathrm{acc})
\cdot \prod_k \mathrm{chiAt}(\text{factor}_k)
= \mathrm{chiAt}(\text{endpoint})$, provided $q$ does not divide the
endpoint.  This is \textbf{proved}
(\lean{EM/Stochastic/RandomTwoPointMC.lean}{chiAt\_multiplicativity\_proved}
{\color{leanink}$\checkmark$})
by induction on~$N$ with backward non-divisibility propagation:
if $q \nmid P_N$, then $q \nmid P_k$ for all $k \leq N$
(\lean{EM/Stochastic/RandomTwoPointMC.lean}{not\_dvd\_epsWalkProdFrom\_of\_le}
{\color{leanink}$\checkmark$}),
so every factor is coprime to~$q$ and $\mathrm{chiAt}$ is
multiplicative.

A full Fourier counting infrastructure is built on top: character
orthogonality on $(\ZZ/q\ZZ)^\times$, death/survival partition
(\lean{EM/Stochastic/RandomTwoPointMCB.lean}{survival\_plus\_death}
{\color{leanink}$\checkmark$}:
$\mathrm{survivalCount} + \mathrm{deathCount} = 2^N$),
the Fourier indicator formula
$|G| \cdot \mathrm{pathCount}(a)
= \sum_\chi \overline{\chi(a)} \cdot U(\chi, N)$,
and norm bounds for the unit endpoint character sum.
The chain is:
\[
  \mathrm{TCA}
  \;\xRightarrow{\checkmark}\;
  \text{vanishing pathCharSum}
  \;\xRightarrow{\text{ChiAtMult}\ \checkmark}\;
  \text{Fourier counting}
  \;\xRightarrow{?}\;
  \text{RandomTwoPointMC}.
\]
The remaining gap is \defname{PathSurvival}
(\lean{EM/Stochastic/RandomTwoPointMCB.lean}{PathSurvival}):
the ratio $\mathrm{survivalCount}(N) / \mathrm{deathCount}(N)$
grows without bound.  Without positive survival at the same~$N$
where pathCharSum is small, the Fourier inversion gives no
information about unit classes.

For $q = 3$, the analysis reveals a structural obstruction.
Since $\mathrm{acc} = 2$ and $\minFac(3) = \operatorname{secondMinFac}(3) = 3$,
the walk immediately dies: $3 \mid P_1 = 6$ on \emph{all} paths,
and absorption is permanent.
Thus PathSurvival(3) is false---but
$\mathrm{TreeContractionAtHalf}(3)$ is \emph{also} false (the tree
character sum is identically~$1$ at every depth, because all
post-death factors are $\equiv 1 \pmod{3}$), making
$\mathrm{TreeContractionImpliesRandomMC}(3)$ vacuously true.

For $q = 3$, the \textbf{zero-mean property} gives maximal per-step
contraction
(\lean{EM/Stochastic/RandomTwoPointMC.lean}{mean\_char\_zero\_at\_diverse\_step\_three}
{\color{leanink}$\checkmark$}):
$(1/2)\chi(a) + (1/2)\chi(b) = 0$ whenever $\chi(a) \neq \chi(b)$
for a nontrivial character on $(\ZZ/3\ZZ)^\times$ (order~$2$, so
$\chi$ maps to $\{1, -1\}$).  Although this property is proved,
the immediate death at $q = 3$ means the zero-mean steps are
never reached from the standard starting point.

\paragraph{TCA $+$ PathSurvival $\Rightarrow$ RandomTwoPointMC.}
The conditional implication is now \textbf{proved}
(\lean{EM/Stochastic/RandomTwoPointMCB.lean}{tca\_path\_survival\_implies\_random\_mc\_proved}
{\color{leanink}$\checkmark$}).
The proof~(${\sim}400$~lines) builds the full Fourier counting
infrastructure on $(\ZZ/q\ZZ)^\times$:
\begin{enumerate}[nosep]
\item \textbf{Path sum decomposition}
  (\lean{EM/Stochastic/RandomTwoPointMCB.lean}{path\_sum\_decomp}
  {\color{leanink}$\checkmark$}):
  the unscaled path sum
  $2^N \cdot \mathrm{pathCharSum}$ splits as
  $\mathrm{chiAt}(2)^{-1}\cdot U(\chi,N) + E(\chi,N)$,
  where $U$ sums over surviving paths via ChiAtMultiplicativity and
  $E$ sums over death paths.

\item \textbf{Death error bound}
  (\lean{EM/Stochastic/RandomTwoPointMCB.lean}{death\_error\_norm\_le}
  {\color{leanink}$\checkmark$}):
  $\|E(\chi,N)\| \leq \mathrm{deathCount}(N)$
  (each death term has norm $\leq 1$).

\item \textbf{Unit endpoint bound}
  (\lean{EM/Stochastic/RandomTwoPointMCB.lean}{uecsi\_norm\_le}
  {\color{leanink}$\checkmark$}):
  $\|U(\chi,N)\| \leq 2^N\|\mathrm{pathCharSum}\| + \mathrm{deathCount}$.

\item \textbf{Fourier counting identity}
  (\lean{EM/Stochastic/RandomTwoPointMCB.lean}{fourier\_counting\_identity}
  {\color{leanink}$\checkmark$}):
  $\sum_\chi \overline{\chi(a)} \cdot U(\chi,N)
  = |G| \cdot \mathrm{pathCount}(a,N)$.

\item \textbf{Fourier lower bound}
  (\lean{EM/Stochastic/RandomTwoPointMCB.lean}{fourier\_lower\_bound}
  {\color{leanink}$\checkmark$}):
  $(q{-}1)\cdot\mathrm{pathCount}(a) \geq
  \mathrm{survivalCount} - (q{-}2)\cdot B$
  where $B$ is the uniform bound on nontrivial $\|U(\chi)\|$.
\end{enumerate}
Combining: TCA gives $\|\mathrm{pathCharSum}\| < 1/(2q)$ for
$N \geq N_0$; PathSurvival gives
$\mathrm{survivalCount} > 2q^2 \cdot \mathrm{deathCount}$ for
$N \geq N_1$.  At $N = \max(N_0,N_1)$, the Fourier lower bound yields
$\mathrm{pathCount}(a) > 0$ for every unit~$a$, via the algebraic
identity $2q^3 + 2q^2 + 3q + 2 > 0$ (true for all $q \geq 5$).
For $q = 3$, TCA is vacuously false, so the implication holds trivially.
The self-consistent RandomTwoPointMC thus reduces to exactly two open
hypotheses: TreeContractionAtHalf and PathSurvival.

\paragraph{Ensemble $\varepsilon$-walk.}
The formalization
(\code{EM/Stochastic/EpsilonWalk.lean}, 458~lines, zero \code{sorry})
develops an ensemble version: instead of fixing $\mathrm{acc} = 2$,
parametrize the $\varepsilon$-walk by squarefree starting point~$n$.
The \defname{SigmaCRTPropagationStep} hypothesis (open) asserts that
CRT equidistribution propagates across steps for $\sigma$-walks,
mirroring CRTPropagationStep for deterministic walks.
Under SRE $+$ SigmaCRTPropagationStep, the death density (fraction of
squarefree~$n$ with $q \mid \mathrm{epsWalkProdFrom}(n,\sigma,k)+1$)
converges to~$1/(q{-}1)$ for every fixed~$k$ and~$\sigma$
(\archived{EM/Archive/Stochastic/EpsilonWalkArchive.lean}{sigma\_death\_density\_tendsto}
{\color{leanink}$\checkmark$}).

\paragraph{Deterministic coverage via Cauchy--Davenport.}
An independent, purely group-theoretic route to coverage is formalized in
\code{EM/Stochastic/IteratedProductCoverage.lean} (293~lines, zero \code{sorry}).
Given a sequence $S_0, S_1, \ldots$ of subsets of a finite commutative
group~$G$ with $\mathrm{minOrder}(G) = |G|$ (true when $|G|$ is prime),
the \defname{iterated product} $P_n = P_{n-1} \cdot S_{n-1}$ starting
from $P_0 = \{1\}$ satisfies
\[
  |P_n| \;\geq\; \min\Bigl(|G|,\; 1 + \sum_{k<n}(|S_k| - 1)\Bigr)
\]
by iterated application of the Cauchy--Davenport theorem
(\lean{EM/Stochastic/IteratedProductCoverage.lean}{iteratedMulFinset\_card\_growth}
{\color{leanink}$\checkmark$}).
If every $|S_k| \geq 2$, then after $|G| - 1$ steps the product is
$\mathrm{Finset.univ}$
(\lean{EM/Stochastic/IteratedProductCoverage.lean}{iteratedMulFinset\_eq\_univ}
{\color{leanink}$\checkmark$}).
For $(\ZZ/p\ZZ)^\times$ with $p{-}1$ prime (safe primes: $p = 5, 7, 11, 23, \ldots$),
$\mathrm{minOrder} = p{-}1 = |G|$, and the coverage bound gives diameter
$p - 2$
(\lean{EM/Stochastic/IteratedProductCoverage.lean}{safe\_prime\_coverage}
{\color{leanink}$\checkmark$}).
This complements the probabilistic Fourier approach: if the
$\varepsilon$-walk factor sets have $|S_k| \geq 2$ at every step (the
diversity condition), then Cauchy--Davenport gives deterministic coverage
in finitely many steps---no spectral hypothesis needed.  The gap is
ensuring diversity: when $\minFac \equiv \operatorname{secondMinFac}
\pmod{q}$, the factor set has size~$1$, and the Cauchy--Davenport
step gives no progress.

\paragraph{Dense capture and Baire category.}
A topological perspective on the two-point variant is developed in
\code{EM/Stochastic/DenseCapture.lean} (301~lines, zero \code{sorry}).
The space of binary selection sequences $\sigma : \NN \to \mathrm{Bool}$
carries the product topology (Cantor space), which is compact Hausdorff
and hence a Baire space.  The \defname{captureSet}$(q, \mathrm{acc})$---the
set of selection sequences that capture prime~$q$ starting from
accumulator~$\mathrm{acc}$---is \emph{open}, unconditionally
(\lean{EM/Stochastic/DenseCapture.lean}{captureSet\_isOpen}
{\color{leanink}$\checkmark$}):
it equals $\bigcup_k \{\sigma \mid \mathrm{factor}(k) = q\}$,
and each level set is open because the factor at step~$k$
depends only on $\sigma(0), \ldots, \sigma(k)$ (finitary dependence,
\lean{EM/Stochastic/DenseCapture.lean}{epsWalkFactorFrom\_depends\_on\_prefix}
{\color{leanink}$\checkmark$}).

The sole open hypothesis is \defname{DenseCaptureHypothesis}: for
$\mathrm{acc} \geq 2$, captureSet$(q, \mathrm{acc})$ is \emph{dense}.
Under this hypothesis, the Baire category theorem gives that the
intersection $\bigcap_{q\;\text{prime}} \mathrm{captureSet}(q, \mathrm{acc})$
is \emph{comeager} (residual): a generic selection sequence captures
every prime
(\lean{EM/Stochastic/DenseCapture.lean}{fullCapture\_residual}
{\color{leanink}$\checkmark$}).
In particular, such a sequence exists
(\lean{EM/Stochastic/DenseCapture.lean}{fullCapture\_nonempty}
{\color{leanink}$\checkmark$}).
A tail restart identity
(\lean{EM/Stochastic/DenseCapture.lean}{epsWalkProdFrom\_tail\_restart}
{\color{leanink}$\checkmark$})
shows the walk from step~$K$ onward is a fresh walk from the
accumulator at step~$K$, connecting the density hypothesis to
the ensemble framework.

%%% The Mixed (1-eps)minFac + eps Random section has been subsumed by
%%% §6.2 (Factor Tree), §6.3 (Factor Confinement), §6.4 (Almost All GenMixedMC).

